%%%%%%%%%%%%%%%%
%%% gen 32 %%%
%%%%%%%%%%%%%%%%


\Chapter{32}


\Verse{1}
{
	\hE{וַיַּשְׁכֵּם לָבָן בַּבֹּקֶר וַיְנַשֵּׁק לְבָנָיו וְלִבְנוֹתָיו וַיְבָרֶךְ אֶתְהֶם וַיֵּלֶךְ וַיָּשָׁב לָבָן לִמְקֹמוֹ׃}
}
{
	\eE{
		And Laban got up early in the morning and kissed his
		sons and his daughters and blessed them and went. And
		Laban went back to his place,
	}
}
{
	Laban gets up early and says goodbye to his progeny like Abe.


	Then he goes home.
}


\Verse{2}
{
	\hE{וְיַעֲקֹב הָלַךְ לְדַרְכּוֹ וַיִּפְגְּעוּ בוֹ מַלְאֲכֵי אֱלֹהִים׃}
}
{
	\eE{and Jacob went his way. And angels of {\God} came upon him.}
}
{
	Jacob goes home too.


	Messengers came on him.
}


\Verse{3}
{
	\hE{וַיֹּאמֶר יַעֲקֹב כַּאֲשֶׁר רָאָם מַחֲנֵה אֱלֹהִים זֶה וַיִּקְרָא שֵׁם הַמָּקוֹם הַהוּא מַחֲנָיִם׃}
}
{
	\eE{
		And Jacob said when he saw them, “This is a camp of {\God},”
		and he called that place’s name Mahanaim.
	}
}
{
	Jacob's like “This is God's camp” so he calls it
	Two Camps because monotheism.\footnote{Some manuscripts have ``because foreshadowing'' here.}
}


\Verse{4}
{
	\hJ{וַיִּשְׁלַח יַעֲקֹב מַלְאָכִים לְפָנָיו אֶל עֵשָׂו אָחִיו אַרְצָה שֵׂעִיר שְׂדֵה אֱדוֹם׃}
}
{
	\eJ{
		And Jacob sent messengers ahead of him to
		Esau, his brother, to the land of Seir, the territory of
		Edom,
	}
}
{
	Jacob sends messengers to Esau his brother who's Seir-y in the land of \paleo{אדם}.
}


\Verse{5}
{
	\hJ{וַיְצַו אֹתָם לֵאמֹר כֹּה תֹאמְרוּן לַאדֹנִי לְעֵשָׂו כֹּה אָמַר עַבְדְּךָ יַעֲקֹב עִם לָבָן גַּרְתִּי וָאֵחַר עַד עָתָּה׃}
}
{
	\eJ{
		and commanded them saying, “You shall say this: ‘To my
		lord, to Esau, your servant Jacob said this, “I’ve stayed
		with Laban and delayed until now,
	}
}
{
	and tells them this, “You shall say this ‘To Esau, Jacob said this “I've stayed with Laban etc.,


%	\footnote{
%		I’ve stayed with Laban. Rashi conveys a tradition that this means: “I
%		lived temporarily with him, but I learned nothing from his bad ways.”
%		But Jacob does in fact learn from Laban. The deceiver meets a deceiver.
%		He learns how it feels to receive an injustice, not just how it feels to
%		dish it out. He learns that it hurts. It is a step in his development as
%		a man. And it apparently brings about the honorable way in which he
%		behaves toward Esau when he returns.
%	}
}


\Verse{6}
{
	\hJ{וַיְהִי לִי שׁוֹר וַחֲמוֹר צֹאן וְעֶבֶד וְשִׁפְחָה וָאֶשְׁלְחָה לְהַגִּיד לַאדֹנִי לִמְצֹא חֵן בְּעֵינֶיךָ׃}
}
{
	\eJ{
		and ox and ass and sheep and male and female servant have
		become mine. And I’m sending to tell my lord so as to find
		favor in your eyes.”’”
	}
}
{
	and I got a lot of ass he says down here three levels deep in the
	quote stack. And I'm sending quote level two to say this to you down
	at level three in the hope that one (that's me) will find favor in
	your eyes.”’”
}


\Verse{7}
{
	\hJ{וַיָּשֻׁבוּ הַמַּלְאָכִים אֶל יַעֲקֹב לֵאמֹר בָּאנוּ אֶל אָחִיךָ אֶל עֵשָׂו וְגַם הֹלֵךְ לִקְרָאתְךָ וְאַרְבַּע מֵאוֹת אִישׁ עִמּוֹ׃}
}
{
	\eJ{
		And the messengers came back to Jacob saying, “We came to
		your brother, to Esau, and also he’s coming to you—and
		four hundred men with him.”
	}
}
{
	They come back to Jacob and they're like:
	
	“He's coming. With four hundred guys.”
}


\Verse{8}
{
	\hJ{וַיִּירָא יַעֲקֹב מְאֹד וַיֵּצֶר לוֹ וַיַּחַץ אֶת הָעָם אֲשֶׁר אִתּוֹ וְאֶת הַצֹּאן וְאֶת הַבָּקָר וְהַגְּמַלִּים לִשְׁנֵי מַחֲנוֹת׃}
}
{
	\eJ{
		And Jacob was very afraid, and he had anguish. And he
		divided the people who were with him and the sheep and the
		oxen and the camels into two camps.
	}
}
{
	Jacob's afraid.


	He splits up into two camps.\footnote{Hey that's the name of the place.}
}


\Verse{9}
{
	\hJ{וַיֹּאמֶר אִם יָבוֹא עֵשָׂו אֶל הַמַּחֲנֶה הָאַחַת וְהִכָּהוּ וְהָיָה הַמַּחֲנֶה הַנִּשְׁאָר לִפְלֵיטָה׃}
}
{
	\eJ{
		And he said, “If Esau will come to one camp and strike it,
		then the camp that is left will survive.”
	}
}
{
	And he's like “They can't kill all of us,” Jacob says, accepting a 50\% casualty rate before the fight begins.
}


\Verse{10}
{
	\hJ{וַיֹּאמֶר יַעֲקֹב אֱלֹהֵי אָבִי אַבְרָהָם וֵאלֹהֵי אָבִי יִצְחָק יְהֹוָה הָאֹמֵר אֵלַי שׁוּב לְאַרְצְךָ וּלְמוֹלַדְתְּךָ וְאֵיטִיבָה עִמָּךְ׃}
}
{
	\eJ{
		And Jacob said, “{\God} of my father Abraham and {\God} of my
		father Isaac, YHWH, who said to me, ‘Go back to your land
		and to your birthplace, and I’ll deal well with you,’
	}
}
{
	The he starts messaging his dad and granddad's God saying how he (the
	God) said ‘Go back your land and to your birthplace and I will deal
	with you real good.’
}


\Verse{11}
{
	\hJ{קָטֹנְתִּי מִכֹּל הַחֲסָדִים וּמִכָּל הָאֱמֶת אֲשֶׁר עָשִׂיתָ אֶת עַבְדֶּךָ כִּי בְמַקְלִי עָבַרְתִּי אֶת הַיַּרְדֵּן הַזֶּה וְעַתָּה הָיִיתִי לִשְׁנֵי מַחֲנוֹת׃}
}
{
	\eJ{
		I’m not worthy of all the kindnesses and all the
		faithfulness that you’ve done with your servant, because I
		crossed this Jordan with just my rod, and now I’ve become
		two camps.
	}
}
{
	He's like “I'm not worthy.”
}


\Verse{12}
{
	\hJ{הַצִּילֵנִי נָא מִיַּד אָחִי מִיַּד עֵשָׂו כִּי יָרֵא אָנֹכִי אֹתוֹ פֶּן יָבוֹא וְהִכַּנִי אֵם עַל בָּנִים׃}
}
{
	\eJ{
		Save me from my brother’s hand, from Esau’s hand, because
		I fear him, in case he’ll come and strike me, mother with
		children.
	}
}
{
	So he's like “Save me from my brother.”
	
	Then he clarifies that he means Esau, in case YHWH didn't know.


	\footnote{
		Jacob asks to be ``delivered'' from Esau (\heb{נצל});
		after wrestling the man later in the chapter he says his life was
		``delivered'' by the same root.
	}


}


\Verse{13}
{
	\hJ{וְאַתָּה אָמַרְתָּ הֵיטֵב אֵיטִיב עִמָּךְ וְשַׂמְתִּי אֶת זַרְעֲךָ כְּחוֹל הַיָּם אֲשֶׁר לֹא יִסָּפֵר מֵרֹב׃}
}
{
	\eJ{
		And you’ve said, ‘I’ll do well with you, and I’ll make
		your seed like the sand of the sea, that it won’t be
		countable because of its great number.’”
	}
}
{
	“Remember what you promised about the seed?”
}


\Verse{14}
{
	\hRJE{וַיָּלֶן שָׁם בַּלַּיְלָה הַהוּא וַיִּקַּח מִן הַבָּא בְיָדוֹ מִנְחָה לְעֵשָׂו אָחִיו׃}
}
{
	\eRJE{
		And he spent that night there. And he took an offering for
		Esau, his brother, from what had come into his hand:
	}
}
{
	He falls asleep and stays the night there.


	He gets a present for Esau.
}


\Verse{15}
{
	\hE{עִזִּים מָאתַיִם וּתְיָשִׁים עֶשְׂרִים רְחֵלִים מָאתַיִם וְאֵילִים עֶשְׂרִים׃}
}
{
	\eE{
		two hundred she-goats and twenty he-goats, two hundred
		ewes and twenty rams,
	}
}
{
	It's a large number of goats.
	
	Also RHLs and AILs.
}


\Verse{16}
{
	\hE{גְּמַלִּים מֵינִיקוֹת וּבְנֵיהֶם שְׁלֹשִׁים פָּרוֹת אַרְבָּעִים וּפָרִים עֲשָׂרָה אֲתֹנֹת עֶשְׂרִים וַעְיָרִם עֲשָׂרָה׃}
}
{
	\eE{
		thirty nursing camels and their offspring, forty cows and
		ten bulls, twenty she-asses and ten he-asses.
	}
}
{
	Also thirty humpy milkers, forty regular milkers, ten bulls,
	and assorted ass of both sexes.
}


\Verse{17}
{
	\hE{וַיִּתֵּן בְּיַד עֲבָדָיו עֵדֶר עֵדֶר לְבַדּוֹ וַיֹּאמֶר אֶל עֲבָדָיו עִבְרוּ לְפָנַי וְרֶוַח תָּשִׂימוּ בֵּין עֵדֶר וּבֵין עֵדֶר׃}
}
{
	\eE{
		And he placed them in his servants’ hands, each herd by
		itself, and he said to his servants, “Pass on in front of
		me, and keep a distance between each herd and the next.”
	}
}
{
	He gives all that to his guys, and he's like:
	
	``As your leader, it's important for you to take the lead.''


	He's the protagonist.
}


\Verse{18}
{
	\hE{וַיְצַו אֶת הָרִאשׁוֹן לֵאמֹר כִּי יִפְגָּשְׁךָ עֵשָׂו אָחִי וִשְׁאֵלְךָ לֵאמֹר לְמִי אַתָּה וְאָנָה תֵלֵךְ וּלְמִי אֵלֶּה לְפָנֶיךָ׃}
}
{
	\eE{
		And he commanded the first, saying, “When my brother,
		Esau, meets you and asks you, saying, ‘To whom do you
		belong, and where are you going, and to whom do these in
		front of you belong?’
	}
}
{
	And he's like ``Ok, when Esau asks who owns you...''
}


\Verse{19}
{
	\hE{וְאָמַרְתָּ לְעַבְדְּךָ לְיַעֲקֹב מִנְחָה הִוא שְׁלוּחָה לַאדֹנִי לְעֵשָׂו וְהִנֵּה גַם הוּא אַחֲרֵינוּ׃}
}
{
	\eE{
		then you’ll say, ‘To your servant, to Jacob. It’s an
		offering sent to my lord, to Esau. And here he is behind
		us as well.’”
	}
}
{
	Then say ``To Jacob, it's a present. Jacob would like to clarify
	that he's your servant and you're his lord. He's at the back.''
}


\Verse{20}
{
	\hE{וַיְצַו גַּם אֶת הַשֵּׁנִי גַּם אֶת הַשְּׁלִישִׁי גַּם אֶת כָּל הַהֹלְכִים אַחֲרֵי הָעֲדָרִים לֵאמֹר כַּדָּבָר הַזֶּה תְּדַבְּרוּן אֶל עֵשָׂו בְּמֹצַאֲכֶם אֹתוֹ׃}
}
{
	\eE{
		And he also commanded the second, also the third, also all
		of those who were going behind the herds, saying, “You’ll
		speak this way to Esau when you find him,
	}
}
{
	He tells the other guys to say the same thing,
	presumably in the hopes that Esau will be fatigued
	from the repetition by the time he reaches the back.
}


\Verse{21}
{
	\hE{וַאֲמַרְתֶּם גַּם הִנֵּה עַבְדְּךָ יַעֲקֹב אַחֲרֵינוּ כִּי אָמַר אֲכַפְּרָה }
	\hE{\hlB{פָנָיו}}
	\hE{ בַּמִּנְחָה הַהֹלֶכֶת לְ}
	\hE{\hlB{פָנַי}}
	\hE{ וְאַחֲרֵי כֵן אֶרְאֶה }
	\hE{\hlB{פָנָיו}}
	\hE{ אוּלַי יִשָּׂא }
	\hE{\hlB{פָנָי׃}}
}
{
	\eE{
		and you’ll say, ‘Here is your servant, Jacob, behind us as
		well.’” Because he said, “Let me appease his \hlB{face} with the
		offering that’s going \hlB{in front of me}, and after that I’ll
		see his \hlB{face}; maybe he’ll raise my \hlB{face}.”
	}
}
{
	He prepares to face his brother.\fC{
		The word ``face'' (\heb{פני}) saturates the episode,
		occurring five times in this sentence alone.
		Later, he says he's seen God ``face to face,''
		and finally tells Esau that seeing his
		face is like seeing God's (33:10).
	}
%	\footnote{
%		\hlB{his face} … \hlB{in front of me} … I’ll see \hlB{his face} … \hlB{my face} … \hlB{ahead of him}.
%		The word “\hlB{face}” recurs five times in this line and all through these two
%		chapters (as does the word “raise”). It also is part of the term that is
%		translated as the preposition “\hlB{in front of}” or “\hlB{ahead of}” (in Hebrew,
%		\hlB{lpny}). The repetition conveys the force of this juncture in Jacob’s
%		life. He must \hlB{face} his past. He must \hlB{face} his brother, whom he wronged.
%		And in the middle of the account of his facing his brother will come the
%		account of his most immediate contact with God in his life, his struggle
%		after which he will say, “I’ve seen God \hlB{face-to-face}.” And the two
%		encounters, first with his God and then with his brother, will then be
%		brought together as he says to Esau, “I’ve seen your \hlB{face}—like seeing
%		God’s \hlB{face}!” (33:10). This is also ironic because Jacob left twenty
%		years earlier in the wake of his deception, which worked because his
%		weak-eyed father could not see his \hlB{face}!
%	}
}


\Verse{22}
{
	\hJ{וַתַּעֲבֹר הַמִּנְחָה עַל }
	\hJ{\hlB{פָּנָיו}}
	\hJ{ וְהוּא לָן בַּלַּיְלָה הַהוּא בַּמַּחֲנֶה׃}
}
{
	\eJ{
		And the offering passed \hlB{ahead of him}. And he spent that
		night in the camp.
	}
}
{
	Everyone else goes ahead of him. He stays back.

	So the gifts are headed toward Esau.

	And Esau's headed toward him.

	Un(conscious+derstandab)ly, Jacob's not moving, cuz he's asleep.\footnote{
		Since they're all moving well below Laban speed in this scene,
		classical biblical mechanics is sufficient for the inferences below.
	}
}


\Verse{23}
{
	\hJ{וַיָּקָם בַּלַּיְלָה הוּא וַיִּקַּח אֶת שְׁתֵּי נָשָׁיו וְאֶת שְׁתֵּי שִׁפְחֹתָיו וְאֶת אַחַד עָשָׂר יְלָדָיו וַיַּעֲבֹר אֵת מַעֲבַר יַבֹּק׃}
}
{
	\eJ{
		And he got up in that night and took his two wives and his
		two maids and his eleven boys and crossed the Jabbok ford.
	}
}
{
	He gets up in the night and takes his two wives and two other wives
	and his eleven boys.
	
	
	He needs to fast as possible here,
	which is presumably why he skips Dinah.\fB{Sic?}


	Then he crosses the Jaboc.\footnote{
		Jaboc (more commonly spelled Jaboq, Jabbok, or Jabok; Hebrew: \heb{יַבֹּק})
		is a biblical river whose name translates to ``pouring out,''
		``emptying,'' or ``flowing rapidly.'' Derived from the Hebrew
		root verb baqaq, it translates literally to ``he will empty.''
	}


%	\footnote{
%		his eleven boys. His daughter, Dinah, is not mentioned.
%	}
}


\Verse{24}
{
	\hJ{וַיִּקָּחֵם וַיַּעֲבִרֵם אֶת הַנָּחַל וַיַּעֲבֵר אֶת אֲשֶׁר לוֹ׃}
}
{
	\eJ{
		And he took them and had them cross the wadi,
		and he had everything that was his cross.
	}
}
{
	All those he bore cross.
}


\Verse{25}
{
	\hJ{וַיִּוָּתֵר יַעֲקֹב לְבַדּוֹ וַיֵּאָבֵק אִישׁ עִמּוֹ עַד עֲלוֹת הַשָּׁחַר׃}
}
{
	\eJ{
		And Jacob was left by himself.
		And a man\footnote{Lit. \heb{איש} in the original.}
		wrestled with him until the dawn’s rising.
	}
}
{
	Strap in folks, this one's a doozy.

	Now Jacob's all by himself.

	He wrestles with a man%
	\footnote{Lit. \paleo{אדמ} in the original.}
	until sunrise.%
	\fA{
		``Wrestled'' is \textit{ʾ-b-q} (\heb{אבק}), a verb found only twice in
		the Bible, both in this scene (Gen.~32:25--26). BDB derives it from
		\textit{ʾabaq}, ``dust'': literally something like ``dust oneself up with
		someone,'' hence grapple or wrestle; Strong's memorably glosses it ``to
		bedust.'' The word may also be playing
		on the nearby river Jabbok (\textit{Yabboq}), whose name in turn sounds
		suspiciously like Jacob (\textit{Yaʿaqob}): Jacob \textit{Yaʿaqob}
		\textit{ʾ-b-q}s at the \textit{Yabboq}. So the famous ``wrestling'' is,
		at the lexical level, two men rolling around and raising dust, inside a
		dense little knot of nearly matching sounds.
	}%
	\fB{
		The ordinary noun means fine,
		easily airborne dust or powder, including the dust kicked up by horses 
		(Ezek.~26:10). For context, Ezekiel 25 is about Ammon, Moab, and 
		Edom; Ezekial 24  is about burning bones and lewdness; and Ezekiel 23
		contains the famous line, ``There she lusted 
		after her lovers, whose genitals were like those of donkeys and whose emission 
		was like that of horses.''
	}%
	\fC{
	`Bless'' is \textit{b-r-k} (\heb{ברך}), the same consonantal root as
		\textit{berek}, `knee.'' The root can also mean to kneel or bow, and the
	plural ``knees'' can mean the lap---the space formed by the knees when
	sitting. Thus the semantic cloud around \textit{b-r-k} includes
	\textit{blessing}, \textit{kneeling/bowing}, \textit{knee}, and
	\textit{lap}. The overlap is especially suggestive here: Jacob demands a
	\textit{b-r-k}, a blessing, in a scene already obsessively concerned with
	his thigh, hip, and groin---and the same three consonants also point toward
	kneeling and the lap.
	}
}




%	\footnote{
%		a man wrestled with him. With whom does he wrestle? It says “a man.” But
%		he is able against the man, and later the man names him yir-’l, which is
%		interpreted to mean “struggles with God,” and says to him, “You’ve
%		struggled with God and with people and were able.” And Jacob names the
%		place Peni-El, meaning “face of God,” and says, “I’ve seen God
%		face-to-face.” This all indicates that he has wrestled with God in human
%		form. But the prophet Hosea refers to this moment and says, “He fought
%		with God, and he fought with an angel and was able” (Hos 12:5–6). This
%		is consistent with the description of angels that I gave in the matter
%		of the three people who visit Abraham (18:3). Angels are material
%		expressions of God’s presence (emanations from the Godhead, hypostases).
%		When they are in front of a human, they look like a “man,” like
%		“people.” That is why Hosea can say “He fought with God” and then say in
%		poetic parallel about the same event “he fought with an angel.”
%	}


\Verse{26}
{
	\hJ{וַיַּרְא כִּי לֹא יָכֹל לוֹ וַיִּגַּע בְּכַף יְרֵכוֹ וַתֵּקַע כַּף יֶרֶךְ יַעֲקֹב בְּהֵאָבְקוֹ עִמּוֹ׃}
}
{
	\eJ{
		And he saw that he was not able against him, and he
		touched the inside of his thigh, and the inside of Jacob’s
		thigh was dislocated during his wrestling with him.
	}
}
{
	And he\fB{Who?} saw that he\fB{Ibid!} was not able against him.\fB{Ibidem!}

	And he%
	\fB{I can't even with this.}
	touched the inside of his%
	\fB{Red can't handle this.}
	thigh%
	\fA{
		The usual translation, ``thigh,'' understates the range of the
		term \heb{ירך}. Across its 34 occurrences (32 verses) in the Masoretic Text,
		it means thigh/hip/haunch (Gen 32:26--33; Judg 3:16,21; Ps 45:4),
		groin/loins (Gen 24:2,9; 47:29; Exod 28:42),
		reproductive source---descendants literally ``come out of'' one's
		\heb{ירך} (Gen 46:26; Exod 1:5; Judg 8:30)---and, spatially, the side/flank
		of the tabernacle or altar (Exod 40:22,24; Lev 1:11; Num 3:29,35)
		and base/shaft of the menorah (Exod 25:31; 37:17; Num 8:4).
		Num 5:21--27 likewise uses it for the woman's reproductive/groin region;
		Song 7:2 for the curves of the thighs/hips;
		Jer 31:19 and Ezek 21:17 for striking the thigh in distress;
		Ezek 24:4 for an animal haunch; and Judg 15:8 in the idiom
		\heb{שוק על ירך}. Thus the range---thigh/hip/haunch $\rightarrow$
		groin/loins/reproductive source, and side/flank $\rightarrow$ base---is
		directly attested in Biblical Hebrew, not merely reconstructed from
		cognates. Comparative caution: Phoenician KAI 4's \textit{yrk} is the
		verb ``prolong,'' not this noun; Akkadian likewise has no simple
		\textit{yarku} ``thigh,'' though \textit{warkatu/urkatu} means
		rear/backside/buttocks (while words such as \textit{qablu} can denote
		loins/hips). Consequently Gen 24:2,9 and 47:29's ``place your hand under my
		\heb{ירך}'' can carry much stronger groin/loins/reproductive force than
		``under my thigh.'' Gen 32:26--33's \heb{כף ירך}, typically translated
		as thigh, is expressed with the same word elsewhere used for
		the reproductive loins. The lexical connection between the oath scenes
		and Jacob's injury is therefore considerably more suggestive than the
		English gloss ``thigh'' makes apparent.
	}

	And the inside of Jacob’s thigh is dislocated.%
	\fC{
		In Gen.~32:26, ``dislocated'' is \textit{watt\=eqa`} (\heb{ותקע}),
		from \textit{t-q-`} (\heb{תקע}), a remarkably broad verb meaning
		drive/thrust in, fix or fasten, strike, clap, or blow a horn;
		``dislocate'' here is inferred from context. It is also the verb used
		throughout Josh.~6 for what the seven priests do with their horns as
		Israel circles Jericho for seven days, with the rear guard following the
		Ark. Most suggestively, the same verb occurs in Num.~25:4 for what is to
		be done to the Israelite leaders after the Baal-Peor apostasy---an episode
		associated with a fertility cult and explicitly with sexual relations
		between Israelite men and Moabite women---where its disputed sense may be
		``hang,'' ``impale,'' or ``expose.'' The underlying physical range is thus
		something like forcefully drive, strike, fix, or leave out of position;
		Gen.~32:26's specific sense ``wrench/dislocate'' is possible, but hardly
		lexically inevitable.
	}


}


\Verse{27}
{
	\hJ{וַיֹּאמֶר שַׁלְּחֵנִי כִּי עָלָה הַשָּׁחַר וַיֹּאמֶר לֹא אֲשַׁלֵּחֲךָ כִּי אִם בֵּרַכְתָּנִי׃}
}
{
	\eJ{
		And he said, “Let me go, because the dawn has risen.” And
		he said, “I won’t let you go unless you bless me.”
	}
}
{
	So he's like ``Let me go, the sun is rising.''%

	Not clear why he's not ok with the sun.%
	\fB{
		Gotta be some kind of demon or vampire.
	}

	He\fB{Who?!} says ``I won't let you go until you kneel.''%
	\fC{
		``Bless'' is \textit{b-r-k} (\heb{ברך}), the same consonantal root as
		\textit{berek}, knee.'' The root can also mean to kneel or bow, and the
		plural ``knees'' can mean the lap---the space formed by the knees when
		sitting. Thus the semantic cloud around \textit{b-r-k} includes
		\textit{blessing}, \textit{kneeling/bowing}, \textit{knee}, and
		\textit{lap}. The overlap is especially suggestive here: Jacob demands a
		\textit{b-r-k}, a blessing, in a scene already obsessively concerned with
		his thigh, hip, and groin---and the same three consonants also point toward
		kneeling and the lap.
	}

	\begingroup
		\BookModeTextSizes
		\scriptsize
		\raggedbottom

		\setlength{\tabcolsep}{0.5em}
		\renewcommand{\arraystretch}{1.15}
	
		% Gen 32:26
		\begin{tabular}{@{}p{0.4875\linewidth}@{\hspace{0.025\linewidth}}p{0.4875\linewidth}@{}}
			\begin{minipage}[t]{\linewidth}
				\vspace{0pt}
				\Table[\relax]{lll}{
					\hJM{וַיֹּאמֶר}
						& \trM{wayyōmer}
						& and he said \\
					\hJM{שַׁלְּחֵנִי}
						& \trM{šalləḥēnî}
						& let me go \\
					\hJM{כִּי}
						& \trM{kî}
						& because \\
					\hJM{עָלָה}
						& \trM{ʿālā}
						& has risen \\
					\hJM{הַשָּׁחַר}
						& \trM{haššaḥar}
						& the dawn \\
				}
			\end{minipage}
			&
			\begin{minipage}[t]{\linewidth}
				\vspace{0pt}
				\Table[\relax]{lll}{
					\hJM{וַיֹּאמֶר}
						& \trM{wayyōmer}
						& and he said \\
					\hJM{לֹא}
						& \trM{lōʾ}
						& not \\
					\hJM{אֲשַׁלֵּחֲךָ}
						& \trM{ʾăšallēḥăḵā}
						& I will let you go \\
					\hJM{כִּי}
						& \trM{kî}
						& except \\
					\hJM{אִם}
						& \trM{ʾim}
						& if \\
					\hJM{בֵּרַכְתָּנִי}
						& \trM{bēraḵtānî}
						& you bless me \\
				}
			\end{minipage}
		\end{tabular}

	\endgroup

%	\footnote{
%		I won’t let you go unless you bless me. Jacob demands a blessing. Is it
%		usual to demand a blessing from someone with whom we have been
%		fighting?! The struggle with God takes such an exertion of one’s
%		strength, and takes such a toll, that the human who does it has to live
%		with the wounds thereafter (Jacob limps after this struggle). And so he
%		is in need of blessing—and has earned it.
%	}
}


\Verse{28}
{
	\hJ{וַיֹּאמֶר אֵלָיו מַה שְּׁמֶךָ וַיֹּאמֶר יַעֲקֹב׃}
}
{
	\eJ{
		And he said to him, “What is your name?” And he said, “Jacob.”
	}
}
{
	Name change incoming.

	\begingroup
		\BookModeTextSizes
		\scriptsize
		\raggedbottom

		\setlength{\tabcolsep}{0.5em}
		\renewcommand{\arraystretch}{1.15}

		% Gen 32:27
		\begin{tabular}{@{}p{0.4875\linewidth}@{\hspace{0.025\linewidth}}p{0.4875\linewidth}@{}}
			\begin{minipage}[t]{\linewidth}
				\vspace{0pt}
				\Table[\relax]{lll}{
					\hJM{וַיֹּאמֶר}
						& \trM{wayyōmer}
						& and he said \\
					\hJM{אֵלָיו}
						& \trM{ʾēlāw}
						& to him \\
					\hJM{מַה}
						& \trM{mah}
						& what \\
					\hJM{שְּׁמֶךָ}
						& \trM{šəmeḵā}
						& your name \\
				}
			\end{minipage}
			&
			\begin{minipage}[t]{\linewidth}
				\vspace{0pt}
				\Table[\relax]{lll}{
					\hJM{וַיֹּאמֶר}
						& \trM{wayyōmer}
						& and he said \\
					\hJM{יַעֲקֹב}
						& \trM{yaʿăqōḇ}
						& Jacob \\
				}
			\end{minipage}
		\end{tabular}
	
		\vspace{1.25\baselineskip}

	\endgroup

}


\Verse{29}
{
	\hJ{וַיֹּאמֶר לֹא יַעֲקֹב יֵאָמֵר עוֹד שִׁמְךָ כִּי אִם יִשְׂרָאֵל כִּי שָׂרִיתָ עִם אֱלֹהִים וְעִם אֲנָשִׁים וַתּוּכָל׃}
}
{
	\eJ{
		And he said, “Your name won’t be said ‘Jacob’ anymore but
		‘Israel,’ because you’ve struggled with {\God} and with
		people and were able.”
	}
}
{
	So he says ``You won't be called\footnote{Lit. ``said'' in the original Eng. and Heb.} Jacob anymore.''

	Now he's\footnote{Lit. he Israel in the original.} Israel.
	
	His etymology is not correct, as far as we know.

	Nobody knows what's going on in this line or why it's here.%
	\fC{
		There are some remarkable Semitic comparanda for the sound sequence
		\textit{r-ʾ-l}. Classical Arabic \textit{raʾl} means a young ostrich,
		with associated imagery of rapid, frightened or wild movement, and
		\textit{Rʾl} is independently attested as a personal name in ancient
		Safaitic. Most strikingly, Ibn Manzur preserves an old poem in which a
		woman propositions a man and touches his penis; in the immediately
		following line the speaker describes what happens to his \textit{raʾl},
		which the lexicographer explains through the behavior of the young
		ostrich. Thus \textit{raʾl} does not lexically mean ``penis,'' but an
		explicitly penile metaphor using the word is independently attested.
		A neighboring but phonologically distinct Arabic root,
		\textit{r-ʿ-l}, means to stab, pierce, or thrust violently, while
		Aramaic \textit{r-ʿ-l} means ``tremble.'' Farther north, the obscure
		form \textit{ʾ-r-ʾ-l} occurs in Hebrew and, remarkably, in the Moabite
		Mesha inscription in connection with Israel; its meaning is disputed,
		with proposals ranging from a cultic object to a human leader.
		No corresponding anatomical \textit{r-ʾ-l} has been securely found in
		Ugaritic or Akkadian. None of this establishes an etymology for
		\textit{Israel}---whose apparent final \textit{r-ʾ-l} crosses the
		ordinary morpheme boundary before \textit{El}---but it supplies an
		unusually rich set of ancient sound-associations for possible
		paronomasia, including, rather inconveniently, an independently
		attested sexual one.
	}


	Suffice it to say, at the end of the day,
	
	Notwithstanding our overall lack of insight into the night fight,

	He won and the man blessed him, \aB{(from+onto) the ground not the other way around.}

	And it's something to do with struggling.

	He's struggled with {\God}\footnote{Lit. \paleo{אל}.}.
	
	And with man\footnote{Lit. \paleo{איש}}.

	That's why he blessed him.
	
	And that's why \paleo{ישראל} is the name we say%
	\fC{
		The root \heb{אמר} is the single most common verb in Biblical Hebrew,
		occurring roughly 5,300 times in the Hebrew Bible. Its basic meaning is
		say, but its very high frequency and use in many different discourse
		contexts give it a considerably wider English translation range. In the
		translation data used here, its renderings, sorted by frequency, include:
		said (2766), saying (862), say (601), says (594), spoke (77), answered
		(50), speak (29), told (25), tell (23), thought (17), spoken (15),
		commanded (14), responded* (7), ordered and promised (6 each), asked and
		\eRed{called} (4 each), commands and resolved (3 each); answer, ask,
		\eRed{call}, continued, decided, declared, follows, gave an order, gave
		the order, gave orders, intend, namely, proposing, said*, speaking, speaks,
		telling, and thinking (2 each); and address, advised, answers, asking,
		asserting, assigned, command, commanded to say, consider, declare*,
		declared*, declares, demonstrates, designate, desired, indeed say,
		informed, intended, intending, meditate, mentioned, name, news, plainly
		says, really thought, repeated*, replied*, requested, saying*, sent word,
		speak to you saying, specifically say, specified, still say, suppose,
		tells, think, utters, and vaunt (1 each).
	} him to this day.

	\begingroup
		\BookModeTextSizes
		\scriptsize
		\raggedbottom

		\setlength{\tabcolsep}{0.5em}
		\renewcommand{\arraystretch}{1.15}

		% Gen 32:28
		\begin{tabular}{@{}p{0.4875\linewidth}@{\hspace{0.025\linewidth}}p{0.4875\linewidth}@{}}
			\begin{minipage}[t]{\linewidth}
				\vspace{0pt}
				\Table[\relax]{lll}{
					\hJM{וַיֹּאמֶר}
						& \trM{wayyōmer}
						& and he said \\
					\hJM{לֹא}
						& \trM{lōʾ}
						& not \\
					\hJM{יַעֲקֹב}
						& \trM{yaʿăqōḇ}
						& Jacob \\
				}
			\end{minipage}
			&
			\begin{minipage}[t]{\linewidth}
				\vspace{0pt}
				\Table[\relax]{lll}{
					\hJM{יֵאָמֵר}
						& \trM{yēʾāmēr}
						& he says \\
					\hJM{עוֹד}
						& \trM{ʿōḏ}
						& again / still / further \\
					\hJM{שִׁמְךָ}
						& \trM{šimḵā}
						& your name \\
					\hJM{כִּי}
						& \trM{kî}
						& but / rather / for / because \\
					\hJM{אִם}
						& \trM{ʾim}
						& if \\
					\hJM{יִשְׂרָאֵל}
						& \trM{yiśrāʾēl}
						& Israel \\
					\hJM{כִּי}
						& \trM{kî}
						& for / because \\
					\hJM{שָׂרִיתָ}
						& \trM{śārîṯā}
						& you contended / struggled \\
					\hJM{עִם}
						& \trM{ʿim}
						& with \\
					\hJM{אֱלֹהִים}
						& \trM{ʾĕlōhîm}
						& gods / God / a god \\
					\hJM{וְעִם}
						& \trM{wəʿim}
						& and with \\
					\hJM{אֲנָשִׁים}
						& \trM{ʾănāšîm}
						& men / people \\
					\hJM{וַתּוּכָל}
						& \trM{wattûḵāl}
						& and you were able / prevailed \\
				}
			\end{minipage}
		\end{tabular}
	
		\vspace{1.25\baselineskip}
	
	\endgroup

%	\footnote{
%		Israel. There is little character development in Adam, Eve, Noah,
%		Abraham, Sarah, Isaac, or Rebekah, all of whom remain basically constant
%		figures through the stories about them. But Jacob changes, and the
%		matter of deception is intimately related to that development. As Esau
%		points out, Jacob’s very name connotes deception: to catch. And Jacob
%		starts out as a manipulator. But Jacob is changed after his experiences
%		in Mesopotamia. He has been the deceiver and the deceived. He has hurt
%		and been hurt. He is now a husband and a father, a man who has struggled
%		and prospered. For the rest of the story he is no longer pictured as a
%		man of action but, more often, as a relatively passive man, like his
%		father Isaac, seeking to appease his brother, avoiding strife and risk.
%		And precisely at the juncture that marks this change in Jacob’s
%		character he has his encounter with God at Penuel. God blesses him in a
%		remarkable etiology/etymology: “Your name shall no longer be called
%		Jacob, but Israel [yir-’l, understood here to mean ‘he struggles with
%		God’], because you’ve struggled with God and with people and were able.
%		Is this divine encounter the signpost of the change in Jacob’s
%		character, or the cause? Either way, as his character changes, and he
%		ceases to be the deceiver, just then he sheds the name Jacob (the one
%		who catches) and becomes instead Israel (the one who struggles with
%		God).
%	}
}


\Verse{30}
{
	\hJ{וַ}
	\hJ{\hlB{יִּשְׁאַל}}
	\hJ{יַעֲקֹב וַיֹּאמֶר הַגִּידָה נָּא שְׁמֶךָ וַיֹּאמֶר לָמָּה זֶּה תִּשְׁאַל לִשְׁמִי וַיְבָרֶךְ אֹתוֹ שָׁם׃}
}
{
	\eJ{
		And Jacob \hlB{asked}. And he said, “Tell your name.” And he
		said, “Why is this that you ask my name?” And he blessed
		him there.
	}
}
{
	So Jacob asks.\fC{
		The verb here is \heb{שאל} (\textit{š-ʾ-l}), conventionally ``ask,'' but
		its semantic range is much wider. Biblical Hebrew uses the root for
		\textit{ask, question, request, petition, seek, desire, demand, beg}.
		One can \heb{שאל} for information, for an object or favor,
		for permission, for someone's welfare, or for an answer, favor, or
		decision from a deity or oracle. The intensity is therefore supplied
		largely by context.	
		
		The same semantic center is ancient and widespread in Semitic. Ugaritic
		\textit{šʾl} means to ask or inquire; Akkadian
		\textit{šaʾālu}/\textit{šâlu} covers asking, questioning, inquiring,
		consulting, asking for, and related petitionary uses; Aramaic and Syriac
		\textit{šʾl} include asking, questioning, consulting, requesting,
		entreating, and demanding; and cognates occur more widely in Semitic.
		The common idea is not specifically ``asking a question,'' but
		\textbf{seeking something from another by asking}---information, an
		answer, an object, permission, assistance, or compliance.

		The grammar of this verse is consequently worth preserving exactly.
		It does \emph{not} say ``Jacob asked him.'' It says
		\textbf{``and Jacob \heb{שאל}, and he said''}; the object of
		\heb{שאל} is unstated. The following \textit{``Tell your name''}
		makes a demand concerning the saying of some name of someone, but
		``him'' is supplied by translators rather than expressed
		by the Hebrew. Nor does the following \textit{``and he said''} identify
		its subject by name: Hebrew simply continues with another masculine
		singular ``he.''
	
		This matters especially in a passage whose masculine pronouns are already
		remarkably slippery. The lexical range permits the first clause to be
		heard in a variety of ways, without specifying in the verb
		itself exactly what was sought or from whom. Only the ensuing speech
		provides the contextual clue. The response then repeats the same root:
		\heb{שאל}, conjugated here as \heb{ישאל}, which, curiously, is one
		letter away from \heb{ישראל} (Israel).
	}\fB{
		Lit. ``Jacob begs'' in the original text.%
	}

	He says ``What's your \paleo{שמ}?''

	He says ``Why ask\footnote{Ibid.} this \paleo{שמ}?''

	He blesses\footnote{
		Lit. \paleo{ברכ} in the original,
		a verb meaning ``to kneel, bend the knee''
		in the standard Qal form. The form
		here is the Piel form, the intensified
		form of the same verb, commonly
		translated as ``bless.''
	} him \paleo{שמ}.

	\begingroup
		\BookModeTextSizes
		\scriptsize
		\raggedbottom

		\setlength{\tabcolsep}{0.5em}
		\renewcommand{\arraystretch}{1.15}
	
		% Gen 32:29
		\begin{tabular}{@{}p{0.4875\linewidth}@{\hspace{0.025\linewidth}}p{0.4875\linewidth}@{}}
			\begin{minipage}[t]{\linewidth}
				\vspace{0pt}
				\Table[\relax]{lll}{
					\hJM{וַיֹּאמֶר}
						& \trM{wayyōmer}
						& and he said \\
					\hJM{הַגִּידָה}
						& \trM{haggîḏā}
						& tell \\
					\hJM{נָּא}
						& \trM{nnā}
						& please \\
					\hJM{שְׁמֶךָ}
						& \trM{šəmeḵā}
						& your name \\
				}
			\end{minipage}
			&
			\begin{minipage}[t]{\linewidth}
				\vspace{0pt}
				\Table[\relax]{lll}{
					\hJM{וַיֹּאמֶר}
						& \trM{wayyōmer}
						& and he said \\
					\hJM{לָמָּה}
						& \trM{lāmmā}
						& why \\
					\hJM{זֶּה}
						& \trM{zeh}
						& this \\
					\hJM{תִּשְׁאַל}
						& \trM{tišʾal}
						& you ask / inquire \\
					\hJM{לִשְׁמִי}
						& \trM{lišmî}
						& for my name \\
					\hJM{וַיְבָרֶךְ}
						& \trM{wayḇāreḵ}
						& and he blessed \\
					\hJM{אֹתוֹ}
						& \trM{ʾōṯô}
						& him \\
					\hJM{שָׁם}
						& \trM{šām}
						& there \\
				}
			\end{minipage}
		\end{tabular}
	
	\endgroup

}


\Verse{31}
{
	\hJ{וַיִּקְרָא יַעֲקֹב שֵׁם הַמָּקוֹם פְּנִיאֵל כִּי רָאִיתִי אֱלֹהִים פָּנִים אֶל פָּנִים וַתִּנָּצֵל נַפְשִׁי׃}
}
{
	\eJ{
		And Jacob called the place’s name Peni-El, “because I’ve
		seen {\God} face-to-face, and my life has been delivered.”
	}
}
{
	\aB{Deliverance!}%
	\fC{%
		The root \heb{נ־צ־ל} has the basic sense of \textbf{taking away, stripping
		away, or pulling something out from somewhere}. In the Niphal, as here,
		it means \textbf{to be delivered, rescued, saved, or to survive}---literally
		something like being \textit{pulled out} of danger. The same root develops
		quite differently in the Piel: \textbf{to make use of, exploit, or abuse}
		someone or something, presumably from the idea of \textit{taking something
		out of} it. Thus ``my life has been delivered'' can retain a concrete
		undertone: Jacob's life has been \textbf{pulled out intact} from his
		face-to-face encounter with {\God}.
	}
}


\Verse{32}
{
	\hJ{וַיִּזְרַח לוֹ הַשֶּׁמֶשׁ כַּאֲשֶׁר עָבַר אֶת פְּנוּאֵל וְהוּא צֹלֵעַ עַל יְרֵכוֹ׃}
}
{
	\eJ{
		And the sun rose on him as he passed Penuel, and he was
		faltering on his thigh.
	}
}
{
	Sun rises as he passes Face of God.

	His parts are a bit off as he heads toward the next section.%
	\footnote{
		Lit. vers secς\heb{כ}ʃиηнne\heb{ה}\heb{ח}x\paleo{t}t+... in the original.
	}

	Note: the text seems a bit blessed%
	\footnote{Lit. blessé, from the French \textsl{blesser}, meaning hurt, wound, injure, or maim.}
	here.
}


\Verse{33}
{
	\hE{עַל כֵּן לֹא יֹאכְלוּ בְנֵי יִשְׂרָאֵל אֶת גִּיד הַנָּשֶׁה אֲשֶׁר עַל כַּף הַיָּרֵךְ עַד הַיּוֹם הַזֶּה כִּי נָגַע בְּכַף יֶרֶךְ יַעֲקֹב בְּגִיד הַנָּשֶׁה׃}
}
{
	\eE{
		On account of this the children of Israel to this day will
		not eat the tendon of the vein that is on the inside of
		the thigh, because he touched the inside of Jacob’s thigh,
		the tendon of the vein.
	}
}
{
	End of class folks, great job all, no further questions.%
	\fC{
		The singular \textit{g\^{i}\d{d}} (``sinew/tendon'') here is effectively a
		hapax in this anatomical construction: outside Gen.~32:32 the noun is rare,
		appearing only in Isa.~48:4 (an ``iron sinew'') and in the plural in
		Job~10:11; 40:17 and Ezek.~37:6,8. The surrounding vocabulary does little
		to make ``ordinary thigh tendon'' secure. \textit{Kaf} denotes a palm,
		hollow, sole, or other concave region; Ezekiel's plural ``sinews'' occurs
		beside bones and flesh in deliberately graphic bodily imagery; and
		Job~40:17 is especially suggestive: immediately after the creature's
		\textit{zanab} (``tail''), its \textit{g\^{i}d\^{e}} are said to be tightly
		interwoven around its dual \textit{pa\d{h}ad}. The latter notoriously
		uncertain word was understood by the KJV/Strong's tradition as ``stones,''
		i.e.\ testicles, while many modern versions euphemize it as ``thighs''; some
		translations accordingly render the whole phrase ``sinews of his male
		organs.'' Relatedly, \textit{pa\d{h}ad} is also the ordinary word
		``fear/dread,'' and
		Genesis itself gives it an extraordinary prominence only one chapter
		earlier: in Gen.~31:42,53 Jacob calls his father's deity the
		``Fear of Isaac'' and then swears ``by the Fear of his father Isaac.''
		The anatomical \textit{pa\d{h}ad} of Job~40:17 and the divine title
		\textit{Pa\d{h}ad Isaac} need not be the same lexical sense; nevertheless,
		the juxtaposition is striking in this particular narrative. Jacob swears by
		Isaac's \textit{pa\d{h}ad}; Abraham's oath had required Jacob's future
		father-in-law's predecessor to put his hand ``under my \textit{yarek}'';
		and shortly afterward Jacob's own \textit{yarek} is struck at its
		\textit{kaf} and associated with a unique singular \textit{gid}.
		Thus Isaiah's ``iron sinew,'' Ezekiel's anatomical usage, Job's
		tail + paired \textit{pa\d{h}ad} + tightly woven sinews, Genesis's
		\textit{kaf} + \textit{yarek} (independently capable of denoting the
		reproductive loins), and Genesis's peculiar ``Fear of Isaac'' terminology
		are not individually decisive. Together they provide several independent
		reasons to suspect that ``sinew of the thigh/hip''---and perhaps even the
		earlier ``put your hand under my thigh''---is considerably more anatomically
		polite than the Hebrew. There are, in short, an awkward number of roads in
		this vocabulary leading toward the family jewels.
	}


}
